% !PRINTED 353
\weberpage{353}

Die letzte Gleichung ergiebt die Unterscheidung der beiden Fälle, in denen die Gleichung c) zwei oder keine positive Wurzel hat.

Man erhält nämlich für das Maximum von $\sin\Theta^{2m}\cos\Theta^{2n}$ nach den Regeln der Differentialrechnung
\[
\frac{m^m n^n}{(m+n)^{m+n}},
\]
das für $\mathrm{tg}\,\Theta^2 = m:n$ erreicht wird. Wenn also $\lambda$ unter dieser Grenze liegt, so haben wir in c) zwei reelle Wurzeln, sonst keine. In den Formeln (2) ist $\lambda$ eine gegebene Grösse, und man hat nun aus den Tafeln den Winkel $\Theta$ zu suchen, der diesen Gleichungen genügt. Wenn man noch gar keine Kenntniss über die ungefähre Lage dieses Winkels hat, so ist es zweckmässig, für die erste Annäherung eine etwa von Grad zu Grad fortschreitende, auf zwei Decimalen abgekürzte Tafel zu benutzen, um den so gefundenen Werth mit Hülfe genauerer Tafeln zu verbessern.

\gesperrt{Gauss} benutzt nicht die trigonometrischen Tafeln, sondern die von ihm zuerst eingeführten Additions- und Subtractions-Logarithmen. Wir wollen dies an einem der Fälle in der Kürze zeigen. Die Einzelheiten für die praktische Anwendung der Methode sind in der \gesperrt{Gauss}'schen Abhandlung zu suchen $^1$).

Die Tafeln der Additions- und Subtractions-Logarithmen, wie sie zuerst von \gesperrt{Gauss} eingeführt und berechnet sind, und wie sie sich jetzt auch in den gebräuchlichen Tabellenwerken finden, geben zu drei Zahlen, die grösser als 1 sind:
\[
a, \quad b = 1+\frac{1}{a}, \quad c = 1+a,
\]
die \gesperrt{Brigg'schen} Logarithmen:
\[
A, \quad B, \quad C.
\]
Ist $\Theta$ ein Winkel im \gesperrt{ersten} Octanten, so kann man setzen:
\[
(3)\quad a = \mathrm{cotg}\,\Theta^2, \quad b = \frac{1}{\cos\Theta^2}, \quad c = \frac{1}{\sin\Theta^2}, \quad 0 < \Theta < 45^0,
\]
und wenn $\Theta$ im \gesperrt{zweiten} Octanten liegt:
\[
(4)\quad a = \mathrm{tg}\,\Theta^2, \quad b = \frac{1}{\sin\Theta^2}, \quad c = \frac{1}{\cos\Theta^2}, \quad 45^0 < \Theta < 90^0.
\]

\noindent\rule{4cm}{0.4pt}

$^1$) Beiträge zur Theorie der algebraischen Gleichungen, zweite Abtheilung (1849). \gesperrt{Gauss} Werke, Bd. III, S. 85.
